CHT Discretization, Acoustic Transient Equations, and LANR Network CouplingEach of the 64 primary cooling sectors contains parallel micro-channels discretized along the flow axis $z$ and transverse plane $(x, y)$. The second-order upwind TVD scheme with the minmod flux limiter is applied to convective terms in the Favre-averaged Navier-Stokes equations:\[
\resizebox{\columnwidth}{!}{$\displaystyle \mathbf{F}_{i+1/2} = \frac{1}{2} \left( \mathbf{F}_i + \mathbf{F}_{i+1} \right) - \frac{1}{2} \phi(r_i) \left( \mathbf{F}_{i+1} - \mathbf{F}_i \right)$}
\]
Supercritical fluid properties $\rho(P, T)$, $\mu(P, T)$, and $k(P, T)$ are evaluated at cell faces via bilinear interpolation from high-resolution tabular datasets.Acoustic pressure wave propagation triggered by the 100 Hz, $10\,\mu\text{s}$ heat flux dump is modeled via the 1D transient perturbation equations:\[
\resizebox{\columnwidth}{!}{$\displaystyle \frac{\partial p'}{\partial t} + \rho_0 c_s^2 \frac{\partial u'}{\partial z} = \frac{\beta_p c_s^2}{c_p} q'''(z, t)$}
\]
\[
\resizebox{\columnwidth}{!}{$\displaystyle \frac{\partial u'}{\partial t} + \frac{1}{\rho_0} \frac{\partial p'}{\partial z} = - \frac{f_{\text{Darcy}} u_0}{2 D_h} u'$}
\]
where $q'''(z, t) = \frac{q''(t) P_{\text{perim}}}{A_{\text{flow}}}$ is the volumetric heat source injected into the near-wall layer. Integrating this hyperbolic wave equation across the channel length demonstrates that the fundamental acoustic resonant frequency $f_{\text{acoustic}} = c_s / (2 L) \approx 122.4\text{ Hz}$ is detuned from the $100\text{ Hz}$ pulse driving frequency, preventing acoustic resonance within the cooling manifold.The dual-stage Half-Heusler and Skutterudite network is modeled as a thermal-electrical network containing $N_{\text{couple}} = 4.80 \times 10^5$ individual thermocouple pairs. The local electrical power generated by couple $k$ operating between hot junction $T_{h, k}$ and cold junction $T_{c, k}$ is governed by:\[
\resizebox{\columnwidth}{!}{$\displaystyle P_{e, k} = I_k^2 R_{\text{load}, k} = \left[ \frac{\int_{T_{c, k}}^{T_{h, k}} S(T) \, dT}{R_{\text{int}, k} + R_{\text{load}, k}} \right]^2 R_{\text{load}, k}$}
\]
Under load-matched conditions ($R_{\text{load}, k} = R_{\text{int}, k}$), maximum power transfer is achieved:\[
\resizebox{\columnwidth}{!}{$\displaystyle P_{e, k}^{\text{max}} = \frac{\left( \bar{S}_k \Delta T_k \right)^2}{4 R_{\text{int}, k}}$}
\]
Aggregating all couples across the high-temperature ($600\text{ K} - 900\text{ K}$) and mid-temperature ($300\text{ K} - 600\text{ K}$) stages yields $P_{\text{elec}} = 447.903\text{ MW}_{\text{elec}}$ at $\eta_{\text{TEG}} = 33.804\%$ from an input thermal inventory $\dot{Q}_{\text{TEG, th}} = 1,325.000\text{ MW}_{\text{th}}$.The electrical interface to the 1,800-module LANR starter grid is governed by Kirchhoff's matrix network equations:\[
\resizebox{\columnwidth}{!}{$\displaystyle \mathbf{I}_{\text{grid}} = \mathbf{Y}_{\text{bus}} \mathbf{V}_{\text{bus}}$}
\]
\[
\resizebox{\columnwidth}{!}{$\displaystyle \sum_{m=1}^{1800} I_{\text{LANR}, m} = I_{\text{TEG}} = 3.5832 \times 10^5\text{ A}$}
\]
with an operating bus potential $V_{\text{bus}} = 1.250\text{ kV}$, maintaining load matching and grid voltage regulation across the power plant's operational envelope.